% !TeX root = ../../../main.tex
\section{Symmetric Powers of Schemes and Torsors}
This section reviews the construction of quotients for schemes and torsors under finite group actions, specifically focusing on symmetric powers. 
To ensure these quotients exist as schemes, we utilize the framework of admissible actions from \cite{sga1}. 
Our treatment here closely follows the exposition in \cite{Guignard2018}
The definitions and results presented below are adapted from their work, the difference being that we ask $\Gamma$ to act on the left, to make the later construction of symmetric powers more natural. 
This foundation provides the necessary criteria for admissibility and base change required to define the symmetric powers of a scheme $X$
 and a $G$-torsor $\mathcal{P}$ over $X$.


Let $S$ be a scheme.

\begin{definition}[(\cite{sga1}, V.1.7).]
    \noindent
    \begin{itemize}
        \item Let $T$ be an object of a category $\mathcal{C}$ endowed with a right action of a group $\Gamma$. We say that \textbf{the quotient $T/\Gamma$ exists} in $\mathcal{C}$ if the covariant functor
    \[
    \begin{aligned}
    \mathcal{C} &\to \text{Sets} \\
    U &\mapsto \text{Hom}_{\mathcal{C}}(T, U)^{\Gamma}
    \end{aligned}
    \]
    is representable by an object of $\mathcal{C}$.
    \item Let $T$ be an $S$-scheme. An action of a finite group $\Gamma$ on $T$ is \textbf{admissible} if there exists an affine $\Gamma$-invariant morphism $f : T \to T'$ such that the canonical morphism $\mathcal{O}_{T'} \to f_* \mathcal{O}_T$ induces an isomorphism from $\mathcal{O}_{T'}$ to $(f_* \mathcal{O}_T)^{\Gamma}$.
\end{itemize}
\end{definition}

%% [REVIEW: never cited -- admissibility facts are used only implicitly; consider labeling + citing, or compressing into prose with the SGA1 citation]
\begin{prop}
    The following holds:
    \begin{enumerate}
        \item \textbf{(\cite{sga1} V.1.3)}. Let $T$ be an $S$-scheme endowed with an admissible right action of a finite group $\Gamma$. If $f : T \to T'$ is an affine $\Gamma$-invariant morphism such that the canonical morphism $\mathcal{O}_{T'} \to f_* \mathcal{O}_T$ induces an isomorphism from $\mathcal{O}_{T'}$ to $(f_* \mathcal{O}_T)^{\Gamma}$, then the quotient $T/\Gamma$ exists and is isomorphic to $T'$.
        \item \textbf{(\cite{sga1}, V.1.8)}. Let $T$ be an $S$-scheme endowed with a right action of a finite group $\Gamma$. Then, the action of $\Gamma$ on $T$ is admissible if and only if $T$ is covered by $\Gamma$-invariant affine open subsets.
        \item \textbf{(\cite{sga1}, V.1.9)}. Let $T$ be an $S$-scheme endowed with an admissible right action of a finite group $\Gamma$, and let $S'$ be a flat $S$-scheme. Then, the action of $\Gamma$ on the $S'$-scheme $T \times_S S'$ is admissible, and the canonical morphism
            \[
            (T \times_S S')/\Gamma \to (T/\Gamma) \times_S S'
            \]
            is an isomorphism.
    \end{enumerate}
\end{prop}



Now let $X$ be an $S$-scheme and let $d \geq 0$ be an integer. The group ${S}_d$ of permutations of $\llbracket 1, d \rrbracket$ acts on the right on the $S$-scheme $X^{\times_S d} = X \times_S \dots \times_S X$ by the formula
\[
(x_i)_{i \in \llbracket 1, d \rrbracket} \cdot \sigma = (x_{\sigma(i)})_{i \in \llbracket 1, d \rrbracket}.
\]

\begin{prop}[\cite{Guignard2018} Proposition 2.27]
    If $X$ is a scheme, Zariski locally quasi-projective over $S$, then the right action of the symmetric group $S_d$ on the $d$-fold fiber product 
    $X^{\times_S d}$ is admissible. 
    Consequently, the quotient $\mathrm{Sym}_S^d(X) = X^{\times_S d}/S_d$ exists as a scheme over $S$.
\end{prop}

\begin{definition}
    Under the hypotheses of the proposition above, we define the \textbf{relative symmetric product} of $X$ over $S$ 
    of degree $d$ as the quotient
    $$\mathrm{Sym}_S^d(X) \coloneqq X^{\times_S d}/S_d.$$
    When the base scheme $S$ is clear from the context, we shall denote this quotient by 
    $X^{(d)}$.
\end{definition}

Guignard shows that when $X=\Spec(B)$ and $S=\Spec(A)$ then $\Sym_{S}^d(X)$ is representable by an affine $S$-scheme (See \cite[Remark~2.28]{Guignard2018}).

\begin{prop}[{\cite[Proposition~2.29]{Guignard2018}}]\label{prop:symmetric_power_flat_base_change}
If $X$ is flat and Zariski-locally quasi-projective over $S$, then $\text{Sym}_S^d(X)$ is flat over $S$. Moreover, for any $S$-scheme $S'$, the canonical morphism
\[
\text{Sym}_{S'}^d(X \times_S S') \to \text{Sym}_S^d(X) \times_S S'
\]
is an isomorphism.
\end{prop}

\subsection*{Geometric Properties of Symmetric Powers}

We record the geometric properties of symmetric powers that will be used in
both main parts of the thesis.

\begin{theorem}\label{theorem:int_geoint_implies_productint}
Let $X$ and $Y$ be integral schemes over a field $k$. If $X$ is geometrically
integral, then $X \times_k Y$ is integral. If both $X$ and $Y$ are
geometrically integral, then $X \times_k Y$ is geometrically integral.
\end{theorem}
\begin{proof}
By \stackstag{038K}, a geometrically integral scheme is both geometrically
irreducible and geometrically reduced. Since $X$ is geometrically irreducible,
the projection
\[
X\times_k Y\longrightarrow Y
\]
induces a bijection on irreducible components by \stackstag{038F}.
Because $Y$ is integral, it follows that $X\times_k Y$ is irreducible.
Moreover, $X$ is geometrically reduced and $Y$ is reduced, so
$X\times_k Y$ is reduced by \stackstag{035Z}. Hence it is integral.

If both $X$ and $Y$ are geometrically integral, the same argument applies
after every field extension $k'/k$. Thus
\[
(X\times_kY)_{k'}\cong X_{k'}\times_{k'}Y_{k'}
\]
is integral for every $k'/k$, and therefore $X\times_kY$ is geometrically
integral.
\end{proof}

\begin{prop}\label{prop:relative_symmetric_power_smooth}
Let $f:X\to S$ be a smooth, Zariski-locally quasi-projective morphism of
relative dimension one. Then the relative symmetric power
$X_S^{(d)}=\operatorname{Sym}_S^d(X)$ is smooth over $S$ of relative dimension
$d$.
\end{prop}

\begin{proof}
The case $d=0$ is immediate. By
\Cref{prop:symmetric_power_flat_base_change}, $X_S^{(d)}$ is flat over $S$,
and its formation commutes with arbitrary base change. It is locally of
finite presentation by noetherian approximation and this base-change
compatibility.

For every geometric point $\bar s\to S$, we have
\[
(X_S^{(d)})_{\bar s}\cong(X_{\bar s})^{(d)}.
\]
Since $X_{\bar s}$ is a smooth curve over an algebraically closed field,
its symmetric power is smooth of dimension $d$ by
\cite[Chapter~III, Proposition~3.2]{milneAV}.
The fiberwise criterion for smoothness (\stackstag{01V8}) therefore proves
the assertion.
\end{proof}

\begin{cor}\label{cor:symmetric_power_curve_geometry}
Let $C$ be a smooth, projective, geometrically connected curve over a field
$k$. For every $d\geq 0$, the scheme $C^{(d)}$ is smooth and geometrically
integral. Moreover, every finite product of symmetric powers of $C$ is
geometrically integral.
\end{cor}

\begin{proof}
The smoothness follows from \Cref{prop:relative_symmetric_power_smooth}.
The curve $C$ is geometrically integral by \stackstag{0366}. After any field
extension, $C^d$ is integral and its image under the quotient map
$C^d\to C^{(d)}$ is irreducible. Since $C^{(d)}$ is smooth, it is geometrically
reduced, and hence geometrically integral. The assertion about finite products
now follows inductively from \Cref{theorem:int_geoint_implies_productint}.
\end{proof}

We shall also need the following symmetric-polynomial calculation.

Let $e_1,\dots,e_d$ be the elementary symmetric polynomials in
$x_1,\dots,x_d$, put $y_i=x_i^{-1}$, and write
\[
    \kappa:=k(e_1/e_d,\dots,e_{d-1}/e_d).
\]

\begin{lemma}[Symmetric regularity]
\label{lemma:symmetric_regularity}
If $\Phi\in k[y_1,\dots,y_d]^{S_d}$ and every monomial of $\Phi$ has total
degree strictly less than $d$, then
\[
    \Phi\in k(e_1/e_d,\dots,e_{d-1}/e_d)=\kappa.
\]
\end{lemma}

\begin{proof}
The identity
\[
    e_j(y_1,\dots,y_d)
       =\frac{e_{d-j}(x_1,\dots,x_d)}{e_d(x_1,\dots,x_d)}
\]
puts $e_1(y),\dots,e_{d-1}(y)$ in $\kappa$. By the fundamental theorem of
symmetric polynomials, each homogeneous part of $\Phi$ is a polynomial in
$e_1(y),\dots,e_d(y)$ with their natural degrees. Since every homogeneous
part has degree $<d$, the generator $e_d(y)$ cannot occur. Hence
$\Phi\in\kappa\subset\kappa[[e_d]]$.
\end{proof}


\subsection*{Symmetric Powers of Torsors}
Let $G$ be a finite abelian group, let $\cP$ be a $G$-torsor over an $S$-scheme $X$ in $S_{\text{Ét}}$. 


%% [REVIEW: never cited -- presumably the input for prop:symmetric_power_torsor (Guignard 2.32), but no proof ever invokes it; cite it there or remove]
\begin{prop}[{[\cite{sga1}]}, IX.5.8]
Assume that $\cP$ and $X$ are endowed with right actions from a finite group $\Gamma$ 
such that the morphism $\cP \to X$ is $\Gamma$-equivariant, and that the following properties hold:
\begin{enumerate}
    \item[(a)] The right $\Gamma$-action on $\cP$ commutes with the left $G$-action.
    \item[(b)] The right $\Gamma$-action on $X$ is admissible, and the quotient morphism $X \to X/\Gamma$ is finite.
    \item[(c)] For any geometric point $\bar{x}$ of $X$, the action of the stabilizer $\Gamma_{\bar{x}}$ of $\bar{x}$ in $\Gamma$ on the fiber $\cP_{\bar{x}}$ of $\cP$ at $\bar{x}$ is trivial.
\end{enumerate}
Then the action of $\Gamma$ on $\cP$ is admissible, and $\cP/\Gamma$ is a $G$-torsor over $X/\Gamma$ in $S_{\text{Ét}}$.
\end{prop}

By \Cref{theorem:torsors_are_realized_as_spaces}, $\cP$ is representable by a finite étale $X$-scheme.
For each $i \in  \llbracket 1, d \rrbracket$ let $p_i : X^{\times_S d} \to X$ be the projection on $i$-th factor, and let us consider the $G$-torsor
\[
p_1^{-1}\cP \otimes^G \cdots \otimes^G p_d^{-1}\cP = \KG \backslash \cP^{\times_S d}
\]
over $X^{\times_S d}$, where $\KG \subseteq G^d$ is the kernel of the multiplication morphism $G^d \to G$ (the dependence on $d$ is suppressed in the notation).
Equivalently, this is the contraction of the $G^d$-torsor
$p_1^{-1}\cP\times\cdots\times p_d^{-1}\cP$ along the multiplication map
$m:G^d\to G$.
Under the identification $G^d/\KG\simeq G$ induced by $m$, the residual
$G$-action is given by
\[
g\cdot[(u_1,\dots,u_d)]
    =[(g_1u_1,\dots,g_du_d)],
\]
where $(g_1,\dots,g_d)$ is any lift of $g$, so that
$g_1\cdots g_d=g$. This is independent of the chosen lift because any two
lifts differ by an element of $\KG$.
The object $\KG \backslash \cP^{\times_S d}$ of $S_{\text{Ét}}$ is too representable by an $S$-scheme which is finite étale over $X^{\times_S d}$.
The group ${S}_d$ acts on the right on $\KG \backslash \cP^{\times_S d}$ by the formula
\[
(p_i)_{i \in \llbracket 1, d \rrbracket} \cdot \sigma = (p_{\sigma(i)})_{i \in \llbracket 1, d \rrbracket}.
\]
This action of ${S}_d$ commutes with the left action of $G$ on $\KG \backslash \cP^{\times_S d}$.

\medskip

\begin{prop}[\cite{Guignard2018} Proposition 2.32.]\label{prop:symmetric_power_torsor}
If $X$ is Zariski-locally quasi-projective on $S$, then the right action of ${S}_d$ on $\KG \backslash \cP^{\times_S d}$ is admissible,
so that the quotient $\cP^{[d]}$ of $\KG \backslash \cP^{\times_S d}$ by ${S}_d$ exists as a scheme over $S$.
Moreover, the canonical morphism $\cP^{[d]} \to \mathrm{Sym}_S^d(X)$ is a $G$-torsor, 
and the morphism
\[
p_1^{-1}\cP \otimes^G \cdots \otimes^G p_d^{-1}\cP \to r^{-1}\cP^{[d]}
\]
where $r : X^{\times_S d} \to \mathrm{Sym}_S^d(X)$ is the canonical projection, is an isomorphism of $G$-torsors over $X^{\times_S d}$.
\end{prop}

\begin{definition}[Symmetric powers of rank-one local systems]
\label{definition:symmetric_power_local_system}
In the setting of the preceding proposition, let $\Lambda$ be a finite
commutative ring, let $G=\Lambda^\times$, and let
$\mathcal{F}$ be a rank-one $\Lambda$-local system on $X$. If
$\mathcal{P}=\operatorname{Fr}(\mathcal{F})$, we define
\[
\mathcal{F}^{[d]}:=\mathcal{P}^{[d]}\times^G\Lambda
\]
on $X^{(d)}$. Thus there is a canonical identification
\[
\operatorname{Fr}\bigl(\mathcal{F}^{[d]}\bigr)
\cong \mathcal{P}^{[d]}.
\]
By \Cref{prop:symmetric_power_torsor}, its pullback along
$r:X^{\times_S d}\to X^{(d)}$ satisfies
\[
r^{-1}\mathcal{F}^{[d]}
\cong p_1^{-1}\mathcal{F}\otimes_\Lambda\cdots\otimes_\Lambda
p_d^{-1}\mathcal{F}.
\]
\end{definition}
%the $S_d$-action on $\KG \backslash \cP^{\times_S d}$ is compatible with the $S_d$-action on $X^{\times_S d}$, so that we get a canonical morphism $\cP^{[d]} \to \mathrm{Sym}_S^d(X)$.

% %Note that there is a canonical map $\cP^{\times_S d}/S_d \to \cP^$, and the $S_d$-action on $\KG \backslash \cP^{\times_S d}$ is compatible with the $S_d$-action on $X^{\times_S d}$, so that we get a canonical morphism $\cP^{[d]} \to \mathrm{Sym}_S^d(X)$.
% \textcolor{red}{This entire corollary is not exact, replace with: 
% https://gemini.google.com/app/8d257f0b8e6f5927
% or
% https://gemini.google.com/app/f36ed1280d780bea
% }

\begin{prop}[Semidirect-product description]
\label{prop:symmetric_power_semidirect_quotient}
Let $\Xi_G:=\KG\rtimes S_d$, where $S_d$ acts on $\KG$ by permuting
coordinates. Then $\Xi_G$ acts on $\cP^{\times_S d}$ by
\[
((g_1,\ldots,g_d),\sigma)\cdot(p_1,\ldots,p_d)
=
(g_1p_{\sigma^{-1}(1)},\ldots,g_dp_{\sigma^{-1}(d)}),
\]
and there is a canonical isomorphism
\[
\cP^{[d]}
\cong
\Xi_G\backslash\cP^{\times_S d}.
\]
Moreover, the formation of $\cP^{[d]}$ commutes with flat base change on
$S$. If, in addition, $X$ is flat over $S$, then it commutes with arbitrary
base change; see \cite[Proposition~2.33]{Guignard2018}.
\end{prop}

\begin{proof}
Because $G$ is abelian, the condition $\prod_{i=1}^d g_i=1$ defining
$\KG$ is invariant under permutation of the coordinates. Thus the
$S_d$-action normalizes $\KG$ and defines the semidirect product $\Xi_G$.
By \Cref{prop:symmetric_power_torsor},
\[
\cP^{[d]}=(\KG\backslash\cP^{\times_S d})/S_d.
\]
The universal properties of the two quotients identify this iterated
quotient canonically with $\Xi_G\backslash\cP^{\times_S d}$. Formation of
the torsor quotient $\KG\backslash\cP^{\times_S d}$ commutes with base
change, and the final admissible $S_d$-quotient commutes with flat base
change by \cite[V.1.9]{sga1}. If $X$ is flat over $S$, compatibility with
arbitrary base change is \cite[Proposition~2.33]{Guignard2018}.
\end{proof}

We record one further compatibility of symmetric powers. It will allow us
later to replace a torsor by the connected torsor under its monodromy
subgroup.

\begin{lemma}[Symmetric powers and induced torsors]
\label{lemma:symmetric_power_induced_torsor}
Let $\iota:H\hookrightarrow G$ be an injection of finite abelian groups, and
let $\cQ$ be an $H$-torsor on $U$. There is a canonical isomorphism of
$G$-torsors over $U^{(d)}$
\[
    (\iota_*\cQ)^{[d]}\cong\iota_*(\cQ^{[d]}).
\]
\end{lemma}

\begin{proof}
Let
\[
    j:\cQ\longrightarrow\iota_*\cQ=G\times^H\cQ
\]
be the canonical morphism $(q\mapsto[1,q])$. It is $H$-equivariant, where
$H$ acts on the target through $\iota$. Applying $j$ to each factor gives
\[
    j^{\times d}:\cQ^{\times d}\longrightarrow(\iota_*\cQ)^{\times d},
    \qquad
    (q_1,\dots,q_d)\longmapsto(j(q_1),\dots,j(q_d)).
\]
This map is compatible with permutations of the factors. Moreover, for
$(h_1,\dots,h_d)\in H^d$,
\[
    j^{\times d}(h_1q_1,\dots,h_dq_d)
      =(\iota(h_1)j(q_1),\dots,\iota(h_d)j(q_d)).
\]
Put $K_H:=\ker(H^d\to H)$. If $(h_1,\dots,h_d)\in K_H$, then
\[
    \iota(h_1)\cdots\iota(h_d)
      =\iota(h_1\cdots h_d)=1,
\]
so $(\iota(h_1),\dots,\iota(h_d))\in\KG$. Hence $j^{\times d}$ descends to
a morphism
\[
    \cQ^{[d]}=(K_H\rtimes S_d)\backslash\cQ^{\times d}
       \longrightarrow
    (\KG\rtimes S_d)\backslash(\iota_*\cQ)^{\times d}
       =(\iota_*\cQ)^{[d]}.
\]
The induced map between the residual groups
$H^d/K_H\simeq H$ and $G^d/\KG\simeq G$ is $\iota$. Therefore this morphism
is $H$-equivariant, where $H$ acts on the target by
$h\cdot x:=\iota(h)\cdot x$.
By the universal property of induction, this extends uniquely to a
$G$-equivariant morphism
\[
    \iota_*(\cQ^{[d]})\longrightarrow(\iota_*\cQ)^{[d]}.
\]
Both sides are $G$-torsors over $U^{(d)}$, so
\Cref{prop:torsor_morphism_iso} shows that it is an isomorphism.
\end{proof}




%
% COMPATIBILITY WITH ADDITION
%
\subsection*{Additivity of Symmetric Powers}
For $d\geq 0$, write
\[
X_S^d:=\underbrace{X\times_S\cdots\times_S X}_{d\text{ factors}},
\qquad X_S^0:=S.
\]
The construction of symmetric powers is compatible with a natural addition
law. For any integers $d_1,d_2\geq 0$, concatenating the factors gives a
canonical isomorphism
\[
c_{d_1,d_2}:X_S^{d_1}\times_S X_S^{d_2}
\xrightarrow{\sim}X_S^{d_1+d_2}.
\]
It is equivariant for the inclusion
$S_{d_1}\times S_{d_2}\hookrightarrow S_{d_1+d_2}$ and therefore induces
the canonical \textbf{addition morphism}
\[
+_{d_1,d_2}:X^{(d_1)}\times_S X^{(d_2)}
\longrightarrow X^{(d_1+d_2)}.
\]

\begin{prop}\label{prop:symmetric_powers_on_curves}
    Let $\cP$ be a $G$-torsor over $X$. There is a canonical isomorphism of $G$-torsors over $X^{(d_1)} \times_S X^{(d_2)}$:
    \[
    +_{d_1, d_2}^{-1} \left( \cP^{[d_1+d_2]} \right) \cong r_1^{-1} \cP^{[d_1]} \otimes^G r_2^{-1} \cP^{[d_2]}.
    \]
\end{prop}
\begin{proof}
Let $\pi = r_{d_1} \times r_{d_2}$. By the commutativity of the following diagram
\[
\begin{tikzcd}
X_S^{d_1}\times_S X_S^{d_2} \arrow[r, "c_{d_1,d_2}", "\sim"'] \arrow[d, "\pi"'] & X_S^{d_1+d_2} \arrow[d, "r_{d_1+d_2}"] \\
X^{(d_1)} \times_S X^{(d_2)} \arrow[r, "+_{d_1, d_2}"] & X^{(d_1+d_2)}
\end{tikzcd}
\]
and applying \Cref{prop:symmetric_power_torsor}, we obtain canonical isomorphisms:
\[
\pi^{-1} \left( +_{d_1, d_2}^{-1} \cP^{[d_1+d_2]} \right) \cong r_{d_1+d_2}^{-1} \cP^{[d_1+d_2]} \cong \sideset{}{^G}\bigotimes_{i=1}^{d_1+d_2} p_i^{-1} \cP
\]
and
\[
\pi^{-1} (r_1^{-1} \cP^{[d_1]} \otimes^G r_2^{-1} \cP^{[d_2]}) \cong \left( \sideset{}{^G}\bigotimes_{i=1}^{d_1} p_i^{-1} \cP \right) \otimes^G \left( \sideset{}{^G}\bigotimes_{j=d_1+1}^{d_1+d_2} p_j^{-1} \cP \right).
\]
Since both pullbacks identify with
\[
\sideset{}{^G}\bigotimes_{k=1}^{d_1+d_2} p_k^{-1} \cP
\]
in an $S_{d_1} \times S_{d_2}$-equivariant manner, the isomorphism descends to the quotient $X^{(d_1)} \times_S X^{(d_2)}$ by \Cref{prop:symmetric_power_flat_base_change}.
\end{proof}

